% File: ai-nhap/2026-11-15_10_de_tham_khao_quan_10_de_2_va_quan_11_de_1.tex
% Nguồn: Tuyển tập đề tham khảo môn Toán Tuyển sinh 10 TP.HCM 2025-2026
% Gồm 14 bài: Quận 10 - Đề 2 (7 bài) và Quận 11 - Đề 1 (7 bài)

% ==============================================================================
% PHẦN 1: ĐỀ THAM KHẢO TUYỂN SINH 10 QUẬN 10 – ĐỀ 2 (7 BÀI)
% ==============================================================================

\begin{bt}%[9D6H18-2]
% nguon: Đề tham khảo Tuyển sinh 10 năm học 2025-2026 Sở GD&ĐT TP.HCM – Quận 10 - Đề 2 (Bài 1)
Cho parabol $(P)\colon y = x^2$.
\begin{enumerate}[a)]
    \item Vẽ đồ thị $(P)$ trên hệ trục tọa độ $Oxy$.
    \item Tìm tọa độ những điểm $M$ của $(P)$ có tung độ và hoành độ bằng nhau.
\end{enumerate}
\loigiai{
\begin{enumerate}[a)]
    \item Bảng giá trị tương ứng của hàm số $y = x^2$:
    \begin{center}
    \begin{tabular}{|c|c|c|c|c|c|}
        \hline
        $x$ & $-2$ & $-1$ & $0$ & $1$ & $2$ \\
        \hline
        $y = x^2$ & $4$ & $1$ & $0$ & $1$ & $4$ \\
        \hline
    \end{tabular}
    \end{center}
    Đồ thị $(P)$ là một đường parabol đi qua các điểm $(-2; 4)$, $(-1; 1)$, $(0; 0)$, $(1; 1)$, $(2; 4)$, nhận trục $Oy$ làm trục đối xứng và có bề lõm quay lên trên.
    \item Điểm $M(x; y)$ thuộc $(P)$ có tung độ và hoành độ bằng nhau khi và chỉ khi
    $y = x \iff x^2 = x \iff x^2 - x = 0 \iff x(x - 1) = 0 \iff \left[\begin{aligned} &x = 0 \implies y = 0 \\ &x = 1 \implies y = 1. \end{aligned}\right.$
    Vậy có hai điểm thỏa mãn yêu cầu đề bài là $M_1(0; 0)$ và $M_2(1; 1)$.
\end{enumerate}
}
\end{bt}

\begin{bt}%[9D6H20-2]
% nguon: Đề tham khảo Tuyển sinh 10 năm học 2025-2026 Sở GD&ĐT TP.HCM – Quận 10 - Đề 2 (Bài 2)
Cho phương trình bậc hai $x^2 - 7x - 4 = 0$.
\begin{enumerate}[a)]
    \item Chứng minh phương trình trên luôn có $2$ nghiệm phân biệt $x_1$, $x_2$.
    \item Không giải phương trình, hãy tính giá trị của biểu thức $A = x_1^3 + x_2^3$.
\end{enumerate}
\loigiai{
\begin{enumerate}[a)]
    \item Phương trình $x^2 - 7x - 4 = 0$ có các hệ số $a = 1$, $c = -4$.
    Vì $a \cdot c = 1 \cdot (-4) = -4 < 0$ nên phương trình luôn có hai nghiệm phân biệt trái dấu $x_1$, $x_2$.
    \item Áp dụng định lí Viète cho phương trình $x^2 - 7x - 4 = 0$, ta có:
    $\begin{cases} S = x_1 + x_2 = -\dfrac{b}{a} = 7 \\ P = x_1 x_2 = \dfrac{c}{a} = -4. \end{cases}$
    Biểu thức $A$ được biến đổi theo $S$ và $P$ như sau:
    $A = x_1^3 + x_2^3 = (x_1 + x_2)^3 - 3x_1 x_2(x_1 + x_2) = S^3 - 3PS$.
    Thay $S = 7$ và $P = -4$ vào ta được:
    $A = 7^3 - 3 \cdot (-4) \cdot 7 = 343 + 84 = 427$.
\end{enumerate}
}
\end{bt}

\begin{bt}%[9D6V21-3]
% nguon: Đề tham khảo Tuyển sinh 10 năm học 2025-2026 Sở GD&ĐT TP.HCM – Quận 10 - Đề 2 (Bài 3)
Một chiếc hộp không có nắp được làm từ một mảnh bìa các-tông. Chiếc hộp có đáy là một hình vuông cạnh $x$ (cm) và chiều cao là $h = 10\text{ cm}$.
\begin{enumerate}[a)]
    \item Viết biểu thức biểu diễn diện tích xung quanh $S_{\text{xq}}$ của chiếc hộp theo $x$.
    \item Biết rằng tổng diện tích xung quanh và diện tích đáy của chiếc hộp là $225\text{ cm}^2$. Tính độ dài cạnh đáy của chiếc hộp.
\end{enumerate}
\loigiai{
\begin{enumerate}[a)]
    \item Chiếc hộp có đáy là hình vuông cạnh $x$ (cm), bốn mặt bên là bốn hình chữ nhật có kích thước $x\text{ cm} \times 10\text{ cm}$.
    Do đó diện tích xung quanh của chiếc hộp là:
    $S_{\text{xq}} = 4 \cdot (x \cdot 10) = 40x\text{ (cm}^2\text{)}$.
    \item Diện tích đáy của chiếc hộp hình vuông là:
    $S_{\text{đáy}} = x^2\text{ (cm}^2\text{)}$ (với $x > 0$).
    Theo đề bài, tổng diện tích xung quanh và diện tích đáy bằng $225\text{ cm}^2$, ta có phương trình:
    $x^2 + 40x = 225 \iff x^2 + 40x - 225 = 0$.
    Ta có $\Delta' = 20^2 - 1 \cdot (-225) = 400 + 225 = 625 = 25^2 > 0$.
    Phương trình có hai nghiệm:
    $x_1 = -20 + 25 = 5$ (nhận vì $x > 0$),
    $x_2 = -20 - 25 = -45$ (loại).
    Vậy độ dài cạnh đáy của chiếc hộp là $5\text{ cm}$.
\end{enumerate}
}
\end{bt}

\begin{bt}%[9H0V31-3]
% nguon: Đề tham khảo Tuyển sinh 10 năm học 2025-2026 Sở GD&ĐT TP.HCM – Quận 10 - Đề 2 (Bài 4)
Một cây lăn sơn tường có phần bông lăn hình trụ với đường kính đáy $6\text{ cm}$ và chiều dài trục là $22\text{ cm}$.
\begin{enumerate}[a)]
    \item Tính diện tích mảng tường sơn được khi cây lăn quay trọn vẹn đúng một vòng (kết quả làm tròn đến hàng phần mười của $\text{cm}^2$).
    \item Bác An cần sơn hai bức tường hình chữ nhật có kích thước giống nhau, mỗi bức tường có chiều dài $10\text{ m}$ và chiều rộng $2{,}5\text{ m}$. Biết rằng mỗi cây lăn sơn chỉ sử dụng được tối đa $320$ vòng quay. Hỏi bác An cần chuẩn bị ít nhất bao nhiêu cây lăn sơn để hoàn thành việc sơn cả hai bức tường?
\end{enumerate}
\loigiai{
\begin{enumerate}[a)]
    \item Phần bông lăn hình trụ có bán kính đáy $R = \dfrac{6}{2} = 3\text{ cm}$ và chiều cao $h = 22\text{ cm}$.
    Khi cây lăn quay được đúng một vòng thì diện tích sơn được chính bằng diện tích xung quanh của hình trụ:
    $S_{\text{xq}} = 2\pi R h = 2\pi \cdot 3 \cdot 22 = 132\pi \approx 414{,}7\text{ cm}^2$.
    \item Tổng diện tích của hai bức tường cần sơn là:
    $S_{\text{tường}} = 2 \cdot (10 \cdot 2{,}5) = 50\text{ m}^2 = 500\,000\text{ cm}^2$.
    Mỗi cây lăn quay tối đa được $320$ vòng, diện tích tối đa mà một cây lăn có thể sơn là:
    $S_{\text{cây}} = 320 \cdot 132\pi \approx 132\,700{,}87\text{ cm}^2$.
    Số cây lăn sơn cần thiết ít nhất để sơn xong hai bức tường là:
    $n = \dfrac{S_{\text{tường}}}{S_{\text{cây}}} = \dfrac{500\,000}{132\,700{,}87} \approx 3{,}768$.
    Vì số cây lăn phải là số nguyên dương nên bác An cần chuẩn bị ít nhất $4$ cây lăn sơn.
\end{enumerate}
}
\end{bt}

\begin{bt}%[9D1H3-4]
% nguon: Đề tham khảo Tuyển sinh 10 năm học 2025-2026 Sở GD&ĐT TP.HCM – Quận 10 - Đề 2 (Bài 5)
Hai bạn An và Bình cùng muốn mua một bộ truyện tranh nhưng riêng mỗi bạn thì không đủ tiền.
An nói với Bình: ``Nếu cậu cho tớ mượn $\dfrac{3}{5}$ số tiền cậu có thì tớ sẽ đủ tiền mua bộ truyện''.
Bình nói với An: ``Nếu cậu cho tớ mượn một nửa số tiền của cậu có thì tớ vừa đủ tiền mua bộ truyện''.
Hỏi mỗi bạn có bao nhiêu tiền? Biết rằng số tiền bạn Bình có nhiều hơn số tiền bạn An có là $40\,000$ đồng.
\loigiai{
Gọi số tiền bạn An có là $x$ (đồng) và số tiền bạn Bình có là $y$ (đồng), điều kiện $x > 0, y > 0$.
Vì số tiền bạn Bình có nhiều hơn bạn An là $40\,000$ đồng nên ta có phương trình:
$y - x = 40\,000 \iff -x + y = 40\,000$. \quad (1)
Giá tiền của bộ truyện tranh được tính theo lời An là:
$x + \dfrac{3}{5}y$ (đồng).
Giá tiền của bộ truyện tranh được tính theo lời Bình là:
$y + \dfrac{1}{2}x$ (đồng).
Do giá tiền bộ truyện tranh là không đổi nên:
$x + \dfrac{3}{5}y = y + \dfrac{1}{2}x \iff \dfrac{1}{2}x - \dfrac{2}{5}y = 0 \iff 5x - 4y = 0$. \quad (2)
Từ (1) và (2) ta có hệ phương trình:
$\begin{cases} -x + y = 40\,000 \\ 5x - 4y = 0 \end{cases} \iff \begin{cases} -4x + 4y = 160\,000 \\ 5x - 4y = 0 \end{cases} \iff \begin{cases} x = 160\,000 \\ y = 200\,000 \end{cases}$ (thỏa mãn điều kiện).
Vậy bạn An có $160\,000$ đồng và bạn Bình có $200\,000$ đồng.
}
\end{bt}

\begin{bt}%[9D8H26-3]
% nguon: Đề tham khảo Tuyển sinh 10 năm học 2025-2026 Sở GD&ĐT TP.HCM – Quận 10 - Đề 2 (Bài 6)
Số điểm tốt của các bạn học sinh lớp 7B đạt được trong tuần qua (gồm 5 ngày học từ thứ Hai đến thứ Sáu) lần lượt là: thứ Hai đạt $9$ điểm tốt, thứ Ba đạt $8$ điểm tốt, thứ Tư đạt $7$ điểm tốt, thứ Năm đạt $8$ điểm tốt, thứ Sáu đạt $6$ điểm tốt. Chọn ngẫu nhiên một ngày học trong tuần đó. Biết rằng khả năng chọn mỗi ngày đều như nhau. Tính xác suất của mỗi biến cố sau:
\begin{enumerate}[a)]
    \item Biến cố $A$: ``Vào ngày được chọn các học sinh lớp 7B đạt đúng $7$ điểm tốt''.
    \item Biến cố $B$: ``Vào ngày được chọn các học sinh lớp 7B đạt ít nhất $8$ điểm tốt''.
\end{enumerate}
\loigiai{
Một tuần học gồm $5$ ngày (thứ Hai, thứ Ba, thứ Tư, thứ Năm, thứ Sáu), do đó tổng số kết quả có thể xảy ra khi chọn ngẫu nhiên một ngày là $n(\Omega) = 5$.
\begin{enumerate}[a)]
    \item Trong $5$ ngày, chỉ có duy nhất ngày thứ Tư là lớp đạt đúng $7$ điểm tốt. Do đó số kết quả thuận lợi cho biến cố $A$ là $n(A) = 1$.
    Xác suất của biến cố $A$ là:
    $P(A) = \dfrac{n(A)}{n(\Omega)} = \dfrac{1}{5} = 0{,}2$.
    \item Các ngày học lớp đạt ít nhất $8$ điểm tốt (tức số điểm tốt $\ge 8$) gồm có: thứ Hai ($9$ điểm tốt), thứ Ba ($8$ điểm tốt), thứ Năm ($8$ điểm tốt). Do đó có $3$ kết quả thuận lợi cho biến cố $B$, tức $n(B) = 3$.
    Xác suất của biến cố $B$ là:
    $P(B) = \dfrac{n(B)}{n(\Omega)} = \dfrac{3}{5} = 0{,}6$.
\end{enumerate}
}
\end{bt}

\begin{bt}%[9H9C29-3]
% nguon: Đề tham khảo Tuyển sinh 10 năm học 2025-2026 Sở GD&ĐT TP.HCM – Quận 10 - Đề 2 (Bài 7)
Cho đường tròn $(O; 3\text{ cm})$. Hai điểm $B, C$ thuộc đường tròn $(O)$ sao cho $\widehat{BOC} = 120^\circ$. Tiếp tuyến của đường tròn $(O)$ tại $B$ và $C$ cắt nhau tại điểm $A$.
\begin{enumerate}[a)]
    \item Chứng minh tứ giác $ABOC$ là tứ giác nội tiếp.
    \item Kẻ đường kính $CE$ của đường tròn $(O)$, đoạn thẳng $AE$ cắt $(O)$ tại điểm $D$ ($D$ khác $E$). Chứng minh $AC^2 = AD \cdot AE$.
    \item Tính độ dài đoạn thẳng $BD$.
\end{enumerate}
\loigiai{
\begin{enumerate}[a)]
    \item Vì $AB, AC$ là các tiếp tuyến của $(O)$ tại $B$ và $C$ nên $OB \perp AB$ và $OC \perp AC$, suy ra $\widehat{ABO} = 90^\circ$ và $\widehat{ACO} = 90^\circ$.
    Xét tứ giác $ABOC$, ta có $\widehat{ABO} + \widehat{ACO} = 90^\circ + 90^\circ = 180^\circ$.
    Do đó tứ giác $ABOC$ là tứ giác nội tiếp đường tròn đường kính $AO$.
    \item Vì $CE$ là đường kính của $(O)$ và $D \in (O)$ nên tam giác $CDE$ nội tiếp chắn nửa đường tròn, suy ra $\widehat{CDE} = 90^\circ$, hay $CD \perp AE$ tại $D$.
    Tam giác $ACE$ vuông tại $C$ (do $AC \perp CE$ tại $C$) có đường cao $CD$.
    Áp dụng hệ thức lượng trong tam giác vuông $ACE$, ta có:
    $AC^2 = AD \cdot AE$ (đpcm).
    \item Tính độ dài đoạn thẳng $BD$:
    Vì $AB, AC$ là hai tiếp tuyến cắt nhau tại $A$ nên $OA$ là tia phân giác của $\widehat{BOC}$, do đó $\widehat{AOC} = \dfrac{\widehat{BOC}}{2} = \dfrac{120^\circ}{2} = 60^\circ$.
    Trong tam giác vuông $ACO$ tại $C$:
    $AC = OC \cdot \tan \widehat{AOC} = 3 \cdot \tan 60^\circ = 3\sqrt{3}\text{ cm}$.
    Đường kính $CE = 2R = 2 \cdot 3 = 6\text{ cm}$.
    Áp dụng định lí Pythagore trong tam giác vuông $ACE$:
    $AE = \sqrt{AC^2 + CE^2} = \sqrt{(3\sqrt{3})^2 + 6^2} = \sqrt{27 + 36} = \sqrt{63} = 3\sqrt{7}\text{ cm}$.
    Vì tam giác $OBC$ cân tại $O$ có $OB = OC = 3\text{ cm}$ và $\widehat{BOC} = 120^\circ$, nên:
    $BC = 2 \cdot OB \cdot \sin 60^\circ = 2 \cdot 3 \cdot \dfrac{\sqrt{3}}{2} = 3\sqrt{3}\text{ cm}$.
    Mặt khác tam giác $EBC$ nội tiếp đường tròn đường kính $CE$ nên vuông tại $B$.
    Do đó $EB = \sqrt{CE^2 - BC^2} = \sqrt{6^2 - (3\sqrt{3})^2} = \sqrt{36 - 27} = \sqrt{9} = 3\text{ cm}$.
    Xét $\triangle ABD$ và $\triangle AEB$, ta có:
    Góc $\widehat{BAE}$ chung;
    $\widehat{ABD} = \widehat{AEB}$ (góc tạo bởi tia tiếp tuyến và dây cung và góc nội tiếp cùng chắn cung $BD$).
    Do đó $\triangle ABD \backsim \triangle AEB$ (g.g), suy ra:
    $\dfrac{BD}{EB} = \dfrac{AB}{AE} \implies BD = EB \cdot \dfrac{AB}{AE}$.
    Vì $AB = AC = 3\sqrt{3}\text{ cm}$, thay số vào ta được:
    $BD = 3 \cdot \dfrac{3\sqrt{3}}{3\sqrt{7}} = \dfrac{3\sqrt{3}}{\sqrt{7}} = \dfrac{3\sqrt{21}}{7}\text{ cm}$.
\end{enumerate}
}
\end{bt}


% ==============================================================================
% PHẦN 2: ĐỀ THAM KHẢO TUYỂN SINH 10 QUẬN 11 – ĐỀ 1 (7 BÀI)
% ==============================================================================

\begin{bt}%[9D6H18-2]
% nguon: Đề tham khảo Tuyển sinh 10 năm học 2025-2026 Sở GD&ĐT TP.HCM – Quận 11 - Đề 1 (Bài 1)
Cho hàm số bậc hai $y = x^2$ có đồ thị $(P)$.
\begin{enumerate}[a)]
    \item Vẽ đồ thị $(P)$ trên mặt phẳng tọa độ $Oxy$.
    \item Tìm tọa độ các điểm trên parabol $(P)$ có tung độ bằng $16$.
    \item Tìm tọa độ các điểm trên parabol $(P)$ (khác gốc tọa độ $O$) cách đều hai trục tọa độ.
\end{enumerate}
\loigiai{
\begin{enumerate}[a)]
    \item Bảng giá trị tương ứng của hàm số $y = x^2$:
    \begin{center}
    \begin{tabular}{|c|c|c|c|c|c|}
        \hline
        $x$ & $-2$ & $-1$ & $0$ & $1$ & $2$ \\
        \hline
        $y = x^2$ & $4$ & $1$ & $0$ & $1$ & $4$ \\
        \hline
    \end{tabular}
    \end{center}
    Đồ thị $(P)$ là một parabol có đỉnh $O(0; 0)$, đối xứng qua trục tung $Oy$ và quay bề lõm lên phía trên.
    \item Các điểm thuộc $(P)$ có tung độ bằng $16$ thỏa mãn:
    $y = 16 \iff x^2 = 16 \iff x = \pm 4$.
    Vậy có hai điểm thỏa mãn là $A(4; 16)$ và $B(-4; 16)$.
    \item Khoảng cách từ một điểm $M(x; y) \in (P)$ đến trục hoành $Ox$ là $|y|$, khoảng cách đến trục tung $Oy$ là $|x|$.
    Điểm $M$ cách đều hai trục tọa độ khi và chỉ khi:
    $|y| = |x| \iff |x^2| = |x| \iff x^2 = |x|$.
    Vì điểm $M$ khác gốc tọa độ $O(0; 0)$ nên $|x| > 0$. Chia cả hai vế cho $|x|$, ta được:
    $|x| = 1 \iff x = 1$ hoặc $x = -1$.
    Với $x = 1 \implies y = 1^2 = 1$, ta có điểm $M_1(1; 1)$.
    Với $x = -1 \implies y = (-1)^2 = 1$, ta có điểm $M_2(-1; 1)$.
    Vậy có hai điểm thỏa mãn yêu cầu bài toán là $(1; 1)$ và $(-1; 1)$.
\end{enumerate}
}
\end{bt}

\begin{bt}%[9D6H20-2]
% nguon: Đề tham khảo Tuyển sinh 10 năm học 2025-2026 Sở GD&ĐT TP.HCM – Quận 11 - Đề 1 (Bài 2)
Cho phương trình bậc hai $3x^2 + 5x - 6 = 0$.
\begin{enumerate}[a)]
    \item Chứng minh phương trình trên luôn có hai nghiệm phân biệt $x_1, x_2$.
    \item Không giải phương trình, hãy tính giá trị của biểu thức $A = (3x_1 - 2x_2)(3x_2 - 2x_1)$.
\end{enumerate}
\loigiai{
\begin{enumerate}[a)]
    \item Phương trình $3x^2 + 5x - 6 = 0$ có hệ số $a = 3$, $c = -6$.
    Vì $a \cdot c = 3 \cdot (-6) = -18 < 0$ nên phương trình luôn có hai nghiệm phân biệt trái dấu $x_1, x_2$.
    \item Áp dụng hệ thức Viète cho phương trình $3x^2 + 5x - 6 = 0$, ta có:
    $\begin{cases} S = x_1 + x_2 = -\dfrac{5}{3} \\ P = x_1 x_2 = \dfrac{-6}{3} = -2. \end{cases}$
    Biến đổi biểu thức $A$:
    $\begin{aligned}
    A &= (3x_1 - 2x_2)(3x_2 - 2x_1) \\
    &= 9x_1 x_2 - 6x_1^2 - 6x_2^2 + 4x_1 x_2 \\
    &= 13x_1 x_2 - 6(x_1^2 + x_2^2) \\
    &= 13x_1 x_2 - 6\left[(x_1 + x_2)^2 - 2x_1 x_2\right] \\
    &= 13P - 6(S^2 - 2P) \\
    &= 25P - 6S^2.
    \end{aligned}$
    Thay $S = -\dfrac{5}{3}$ và $P = -2$ vào biểu thức $A$, ta được:
    $A = 25 \cdot (-2) - 6 \cdot \left(-\dfrac{5}{3}\right)^2 = -50 - 6 \cdot \dfrac{25}{9} = -50 - \dfrac{50}{3} = -\dfrac{200}{3}$.
\end{enumerate}
}
\end{bt}

\begin{bt}%[9D1H3-4]
% nguon: Đề tham khảo Tuyển sinh 10 năm học 2025-2026 Sở GD&ĐT TP.HCM – Quận 11 - Đề 1 (Bài 3)
Nguyễn Thị Ánh Viên -- kình ngư xuất sắc của bơi lội Việt Nam -- sinh ngày 9 tháng 11 năm $\overline{19ab}$ (với $a, b \in \{0; 1; \dots; 9\}$).
\begin{enumerate}[a)]
    \item Hãy viết biểu thức biểu thị số tuổi của Ánh Viên vào đúng ngày sinh nhật của cô trong năm 2021 theo hai chữ số $a$ và $b$.
    \item Em hãy xác định chính xác năm sinh của Ánh Viên. Biết rằng đến đúng ngày sinh nhật của mình trong năm 2021, cô nhận thấy số tuổi của mình bằng tổng các chữ số trong năm sinh của cô.
\end{enumerate}
\loigiai{
\begin{enumerate}[a)]
    \item Năm sinh của Ánh Viên là $\overline{19ab} = 1900 + 10a + b$.
    Vào đúng ngày sinh nhật 9 tháng 11 năm 2021, số tuổi của cô được biểu thị bởi:
    $T = 2021 - (1900 + 10a + b) = 121 - 10a - b$.
    \item Tổng các chữ số trong năm sinh của Ánh Viên là:
    $S = 1 + 9 + a + b = 10 + a + b$.
    Theo giả thiết, đến ngày sinh nhật năm 2021, số tuổi bằng tổng các chữ số của năm sinh, ta có phương trình:
    $121 - 10a - b = 10 + a + b \iff 11a + 2b = 111$.
    Vì $b \le 9$ nên $2b \le 18 \implies 11a = 111 - 2b \ge 111 - 18 = 93 \implies a \ge \dfrac{93}{11} \approx 8{,}45$.
    Do $a$ là chữ số nên $a = 9$.
    Với $a = 9$, ta có:
    $11 \cdot 9 + 2b = 111 \iff 99 + 2b = 111 \iff 2b = 12 \iff b = 6$ (thỏa mãn $b$ là chữ số).
    Vậy năm sinh của Ánh Viên là năm 1996.
\end{enumerate}
}
\end{bt}

\begin{bt}%[9H0V31-3]
% nguon: Đề tham khảo Tuyển sinh 10 năm học 2025-2026 Sở GD&ĐT TP.HCM – Quận 11 - Đề 1 (Bài 4)
Một chiếc ly đựng nước có dạng hình trụ với chiều cao $15\text{ cm}$, đường kính đáy là $5\text{ cm}$. Mực nước tinh khiết trong ly đang cao $10\text{ cm}$ (bỏ qua bề dày của đáy và thành ly). Lấy $\pi \approx 3{,}14$.
\begin{enumerate}[a)]
    \item Lượng nước hiện đang chứa trong ly là bao nhiêu centimét khối? (Kết quả làm tròn đến hàng phần trăm).
    \item Người ta thả vào ly nước $5$ viên bi hình cầu có cùng thể tích, đồng chất và chìm ngập hoàn toàn trong nước, làm nước trong ly dâng lên vừa đúng bằng miệng ly. Hỏi bán kính của mỗi viên bi là bao nhiêu centimét? (Kết quả làm tròn đến hàng phần trăm).
\end{enumerate}
\loigiai{
Bán kính đáy của chiếc ly hình trụ là $R = \dfrac{5}{2} = 2{,}5\text{ cm}$.
\begin{enumerate}[a)]
    \item Thể tích lượng nước chứa trong ly là:
    $V_{\text{nước}} = \pi R^2 h_{\text{nước}} \approx 3{,}14 \cdot (2{,}5)^2 \cdot 10 = 3{,}14 \cdot 6{,}25 \cdot 10 = 196{,}25\text{ cm}^3$.
    \item Khi thả $5$ viên bi ngập hoàn toàn trong nước, mực nước dâng từ độ cao $10\text{ cm}$ lên miệng ly có độ cao $15\text{ cm}$.
    Chiều cao cột nước dâng thêm là:
    $h_{\text{dâng}} = 15 - 10 = 5\text{ cm}$.
    Tổng thể tích của $5$ viên bi bằng thể tích của phần nước dâng lên:
    $V_{5\text{ bi}} = \pi R^2 h_{\text{dâng}} \approx 3{,}14 \cdot (2{,}5)^2 \cdot 5 = 98{,}125\text{ cm}^3$.
    Thể tích của một viên bi là:
    $V_{\text{bi}} = \dfrac{98{,}125}{5} = 19{,}625\text{ cm}^3$.
    Mặt khác, thể tích mỗi viên bi hình cầu bán kính $r$ được tính bởi công thức:
    $V_{\text{bi}} = \dfrac{4}{3}\pi r^3 \implies \dfrac{4}{3} \cdot 3{,}14 \cdot r^3 = 19{,}625$.
    Suy ra:
    $r^3 = \dfrac{19{,}625 \cdot 3}{4 \cdot 3{,}14} = \dfrac{58{,}875}{12{,}56} = 4{,}6875 \implies r = \sqrt[3]{4{,}6875} \approx 1{,}67\text{ cm}$.
    Vậy bán kính của mỗi viên bi xấp xỉ $1{,}67\text{ cm}$.
\end{enumerate}
}
\end{bt}

\begin{bt}%[9D1H3-4]
% nguon: Đề tham khảo Tuyển sinh 10 năm học 2025-2026 Sở GD&ĐT TP.HCM – Quận 11 - Đề 1 (Bài 5)
Một trường học tổ chức cho $140$ người gồm giáo viên và học sinh tham gia chuyến tham quan du lịch sinh thái. Giá vé vào cổng niêm yết là $100\,000$ đồng mỗi người. Do được ưu đãi, giá vé vào cổng dành cho học sinh được giảm $10\%$ so với giá niêm yết, do đó tổng số tiền vé nhà trường đã thanh toán là $12\,720\,000$ đồng. Hỏi có bao nhiêu giáo viên và bao nhiêu học sinh đã tham gia chuyến đi?
\loigiai{
Gọi số giáo viên tham gia chuyến đi là $x$ người và số học sinh tham gia là $y$ người ($x, y \in \mathbb{N}^*$, $x, y < 140$).
Vì tổng số người tham gia là $140$ người nên ta có phương trình:
$x + y = 140$. \quad (1)
Giá vé vào cổng của mỗi giáo viên là $100\,000$ đồng.
Giá vé vào cổng của mỗi học sinh sau khi được giảm $10\%$ là:
$100\,000 \cdot (1 - 0{,}10) = 90\,000$ đồng.
Tổng số tiền vé nhà trường đã trả là $12\,720\,000$ đồng nên ta có phương trình:
$100\,000x + 90\,000y = 12\,720\,000 \iff 10x + 9y = 1272$. \quad (2)
Từ (1) và (2) ta có hệ phương trình:
$\begin{cases} x + y = 140 \\ 10x + 9y = 1272 \end{cases} \iff \begin{cases} 9x + 9y = 1260 \\ 10x + 9y = 1272 \end{cases} \iff \begin{cases} x = 12 \\ y = 140 - 12 = 128 \end{cases}$ (thỏa mãn điều kiện).
Vậy có $12$ giáo viên và $128$ học sinh tham gia chuyến đi.
}
\end{bt}

\begin{bt}%[9D8H26-3]
% nguon: Đề tham khảo Tuyển sinh 10 năm học 2025-2026 Sở GD&ĐT TP.HCM – Quận 11 - Đề 1 (Bài 6)
Bạn An có một ổ khóa số dùng cho xe đạp gồm ba vòng xoay, trên mỗi vòng có các chữ số từ $0$ đến $9$ xếp theo thứ tự tròn. Khóa sẽ phát ra một tiếng tách nhẹ khi quay lên hoặc quay xuống đúng một chữ số trên mỗi vòng (kể cả khi quay từ $0$ sang $9$ hay ngược lại). Hiện tại ba vòng xoay đang ở vị trí $9 - 0 - 4$, trong khi mật mã đúng để mở khóa là $5 - 8 - 7$.
\begin{enumerate}[a)]
    \item Tính số tiếng tách ít nhất bạn An cần quay để mở được ổ khóa.
    \item Một người bạn của An cũng mở được khóa từ vị trí $9 - 0 - 4$ nhưng lại quay theo hướng làm cho số tiếng tách là nhiều nhất. Tính số tiếng tách trung bình cần để mở được ổ khóa (xem như số này bằng trung bình cộng của số tiếng tách ít nhất và số tiếng tách nhiều nhất).
\end{enumerate}
\loigiai{
Mỗi vòng xoay có $10$ chữ số xếp theo vòng tròn. Khi chuyển giữa hai số bất kỳ trên cùng một vòng, có hai hướng quay: một hướng quay ngắn hơn và một hướng quay dài hơn, tổng số bước quay của hai hướng luôn bằng $10$.
\begin{enumerate}[a)]
    \item Xét số tiếng tách ít nhất ở từng vòng xoay:
    \begin{itemize}
        \item Vòng thứ nhất (từ số $9$ sang số $5$):
        Quay lùi $9 \to 8 \to 7 \to 6 \to 5$ mất $4$ tiếng tách; quay tiến qua $0$ mất $10 - 4 = 6$ tiếng tách. Số tiếng tách ít nhất là $4$.
        \item Vòng thứ hai (từ số $0$ sang số $8$):
        Quay lùi qua $9$ mất $2$ tiếng tách ($0 \to 9 \to 8$); quay tiến mất $10 - 2 = 8$ tiếng tách. Số tiếng tách ít nhất là $2$.
        \item Vòng thứ ba (từ số $4$ sang số $7$):
        Quay tiến $4 \to 5 \to 6 \to 7$ mất $3$ tiếng tách; quay lùi qua $0$ mất $10 - 3 = 7$ tiếng tách. Số tiếng tách ít nhất là $3$.
    \end{itemize}
    Vậy số tiếng tách ít nhất cần quay để mở khóa là:
    $N_{\min} = 4 + 2 + 3 = 9$ (tiếng tách).
    \item Xét số tiếng tách nhiều nhất ở từng vòng:
    \begin{itemize}
        \item Vòng 1 nhiều nhất là $6$ tiếng tách.
        \item Vòng 2 nhiều nhất là $8$ tiếng tách.
        \item Vòng 3 nhiều nhất là $7$ tiếng tách.
    \end{itemize}
    Số tiếng tách nhiều nhất khi quay mở khóa là:
    $N_{\max} = 6 + 8 + 7 = 21$ (tiếng tách).
    Số tiếng tách trung bình cần để mở được ổ khóa là:
    $N_{\text{tb}} = \dfrac{N_{\min} + N_{\max}}{2} = \dfrac{9 + 21}{2} = 15$ (tiếng tách).
\end{enumerate}
}
\end{bt}

\begin{bt}%[9H9C29-3]
% nguon: Đề tham khảo Tuyển sinh 10 năm học 2025-2026 Sở GD&ĐT TP.HCM – Quận 11 - Đề 1 (Bài 7)
Từ điểm $S$ nằm bên ngoài đường tròn $(O; R)$ vẽ hai tiếp tuyến $SM, SN$ đến đường tròn $(O)$ ($M, N$ là các tiếp điểm và $OS > 2R$).
\begin{enumerate}[a)]
    \item Chứng minh tứ giác $SMON$ là tứ giác nội tiếp.
    \item Gọi $K$ là giao điểm của đoạn thẳng $OS$ với đường tròn $(O)$; $H$ là giao điểm của $SO$ và $MN$. Chứng minh $MH \cdot KS = KH \cdot MS$.
    \item Cho biết $OS = R\sqrt{5}$. Tính bán kính đường tròn nội tiếp của tam giác $SMN$ theo $R$.
\end{enumerate}
\loigiai{
\begin{enumerate}[a)]
    \item Vì $SM, SN$ là hai tiếp tuyến của đường tròn $(O)$ nên $SM \perp OM$ và $SN \perp ON$, suy ra $\widehat{SMO} = 90^\circ$ và $\widehat{SNO} = 90^\circ$.
    Xét tứ giác $SMON$, ta có $\widehat{SMO} + \widehat{SNO} = 90^\circ + 90^\circ = 180^\circ$.
    Do đó tứ giác $SMON$ nội tiếp đường tròn đường kính $SO$.
    \item Theo tính chất hai tiếp tuyến cắt nhau, ta có $SM = SN$ và $OM = ON = R$, suy ra $SO$ là đường trung trực của đoạn thẳng $MN$.
    Do đó $SO \perp MN$ tại $H$, suy ra tam giác $SMH$ vuông tại $H$.
    Xét tam giác $OMK$ có $OM = OK = R$ nên tam giác $OMK$ cân tại $O$, do đó $\widehat{OMK} = \widehat{OKM}$.
    Mặt khác, trong tam giác vuông $SMO$ tại $M$, ta có:
    $\widehat{SMK} + \widehat{KMO} = 90^\circ \implies \widehat{SMK} = 90^\circ - \widehat{KMO}$.
    Trong tam giác vuông $KMH$ tại $H$, ta có:
    $\widehat{KMH} = 90^\circ - \widehat{MKH} = 90^\circ - \widehat{MKO}$.
    Từ đó suy ra $\widehat{SMK} = \widehat{KMH}$.
    Do đó $MK$ là tia phân giác trong của góc $\widehat{SMH}$ trong tam giác $SMH$.
    Áp dụng tính chất đường phân giác trong tam giác $SMH$, ta có:
    $\dfrac{KH}{KS} = \dfrac{MH}{MS} \implies MH \cdot KS = KH \cdot MS$ (đpcm).
    \item Xét tam giác $SMO$ vuông tại $M$ có $OM = R$, $OS = R\sqrt{5}$:
    $SM = \sqrt{OS^2 - OM^2} = \sqrt{(R\sqrt{5})^2 - R^2} = \sqrt{5R^2 - R^2} = 2R$.
    Áp dụng hệ thức lượng trong tam giác vuông $SMO$ có đường cao $MH$:
    $MH \cdot OS = SM \cdot OM \implies MH = \dfrac{SM \cdot OM}{OS} = \dfrac{2R \cdot R}{R\sqrt{5}} = \dfrac{2R}{\sqrt{5}}$.
    $SH = \dfrac{SM^2}{OS} = \dfrac{(2R)^2}{R\sqrt{5}} = \dfrac{4R}{\sqrt{5}}$.
    Đoạn $OK = R$, suy ra $SK = OS - OK = R\sqrt{5} - R = R(\sqrt{5} - 1)$.
    Độ dài đoạn thẳng $KH$ là:
    $KH = SH - SK = \dfrac{4R}{\sqrt{5}} - R(\sqrt{5} - 1) = \dfrac{4R - 5R + R\sqrt{5}}{\sqrt{5}} = \dfrac{R(\sqrt{5} - 1)}{\sqrt{5}}$.
    Vì $SO$ là trục đối xứng của tam giác $SMN$ nên $SK$ là tia phân giác của góc $\widehat{MSN}$.
    Lại có $MK$ là tia phân giác của $\widehat{SMN}$ (do $MK$ là phân giác của $\widehat{SMH}$).
    Suy ra $K$ chính là tâm đường tròn nội tiếp tam giác $SMN$.
    Vì $SO \perp MN$ tại $H$ nên khoảng cách từ $K$ đến cạnh $MN$ chính là đoạn thẳng $KH$.
    Do đó bán kính đường tròn nội tiếp tam giác $SMN$ là:
    $r = KH = \dfrac{R(\sqrt{5} - 1)}{\sqrt{5}} = \dfrac{R(5 - \sqrt{5})}{5}$.
\end{enumerate}
}
\end{bt}
